\documentclass[12pt,review,authoryear,nopreprintline]{elsarticle}
\input{preamble}
\usepackage[margin=1in]{geometry}
\usepackage{setspace}
\makeatletter
\AtBeginDocument{\doublespacing\global\let\@blstr\baselinestretch}
\makeatother
\usepackage{float}
\usepackage[labelfont=bf,labelsep=space]{caption}
\renewcommand{\figurename}{Fig.}
\providecommand{\DOIprefix}{}
\providecommand{\doi}[1]{\url{https://doi.org/#1}}
\begin{document}
\begin{frontmatter}
\title{Private Compatibility and Government Incentives for Standards Harmonization}
\author{Anonymous Author}
\begin{abstract}
When can private adaptation to an incompatible regional standard change governments' incentives to negotiate a wider agreement? We study three countries whose governments choose separate, regional, or common standards before firms decide whether to incur a fixed cost to support an additional standard. In a segmented Cournot model, one outsider's adoption erodes the protected member-market rents of a regional union without reciprocally removing members' disadvantage in the outsider market. With residual per-unit adaptation cost $d$, the member's regional-minus-international welfare difference is $-\mathscr D+R(d,v)$, where $R$ is obtained from the re-solved product market. This yields precisely bounded preference reversal, stronger incentive without reversal, and weaker incentive regions. Strict outsider-only adoption and a government preference reversal coexist on an open set even when adaptation cost relief is incomplete. The political effect is sensitive to the network and demand structure: the initial regional advantage vanishes at zero network effects, under one alternative singleton-network assumption, and in a pre-specified differentiated-demand test under Cournot and Bertrand competition. A full-bypass coalition-stability application distinguishes robust international coordination from blocking-rule-sensitive regional stability. We offer a conditional mechanism, not a general prediction that interoperability promotes international harmonization.
\end{abstract}
\begin{keyword}
private compatibility \sep standards \sep interoperability \sep imperfect competition \sep coalition formation \sep international harmonization
\JEL L13 \sep L15 \sep F15 \sep C71
\end{keyword}
\end{frontmatter}
\makeatletter
\typeout{JICT BODY FONT=\f@size; BASELINE=\the\baselineskip}
\makeatother
\section{Introduction}\label{sec:short-intro}
Standards may be formally incompatible but privately interoperable. A foreign producer can incur a setup cost to support another country's technical interface or certification, reducing a government's ability to exclude the producer without changing the formal membership of a regional agreement. The institutional response is not obvious: private workarounds raise member-market consumer surplus but also erode domestic producers' competitive rents, while the members' reciprocal disadvantage in an outsider market may persist.

We ask when a privately chosen, \emph{one-way} adaptation changes a regional member government's ranking of regional and international standardization. We distinguish the private adoption equilibrium from the formal institutional outcome. In a three-country Cournot environment, the pre-adaptation regional-minus-international welfare gap is $\mathscr E-\mathscr D$: a net member-market advantage less an outsider-market disadvantage. After outsider adaptation it is $-\mathscr D+R(d,v)$, where $R(d,v)$ is the residual member-market rent derived from market equilibrium. Depending on the residual cost, private adaptation can reverse the preference, strengthen the incentive without reversal, or weaken it relative to no adoption. An exact joint witness ensures that the private action generating the political effect really occurs in equilibrium.

Private converters and one-way compatibility are established firm choices \citep{farrellSaloner1992,manentiSomma2008}. \citet{gandalShy2001} show that conversion costs can sustain standardization unions, whereas strong network effects favor mutual recognition. \citet{takaradaEtAl2020} find multilateral agreements in the core for sufficiently local consumption externalities and regional agreements for large transboundary externalities; their fixed-cost extension changes regime payoffs. Here the additional-standard investment is instead chosen after the formal partition, changing the continuation value of the same union. \citet{innocentiMenicucci2021} characterize firms' compatibility choices before price competition, and \citet{dingKoShen2022} study coalition-proof partial compatibility among private platforms. Our equilibrium object is a national government's ranking of formal partitions after firms' adoption decisions. \citet{schmidtSteingress2022} connect costly firm adoption to trade responses to harmonized standards; we ask how the resulting market rents affect incentives to change a formal agreement.

A working paper by Maggi and Mr{\'{a}}zov{\'{a}} (NBER Working Paper 33318, 2024) argues that regulatory harmonization may arise without an agreement and that cooperation may instead preserve efficient diversity. Our contribution is an explicit, \emph{equilibrium-feasible} mapping from private post-agreement adaptation to \emph{national} incentives over formal standards regimes, together with failure boundaries. The coalition-stability result is a full-adaptation application, not a theorem about incomplete-adaptation stability.

\paragraph{Relation to public policy and private standards choice} \citet{fischerSerra2000} and \citet{grossmanMcCalmanStaiger2021} study the costs of serving markets governed by different standards, while \citet{klimenko2009} analyzes public policies and international agreements concerning technical compatibility. \citet{kawabataTakarada2021} examine whether regional standards cooperation impedes wider harmonization. On the firm side, \citet{buccellaEtAl2023} study endogenous compatibility with network effects, and \citet{parentiVannoorenberghe2025} link costly local product tailoring to regulatory blocs. Regulatory product standards can address consumption externalities, whereas compatibility here changes demand and producer cost. A different government objective alone is not our contribution: the model derives the national welfare effect of an outsider's private adoption after a formal partition is selected, including the firm continuation that makes that adoption feasible.

Section~\ref{sec:short-model} sets out the primitives. Sections~\ref{sec:short-adoption} and \ref{sec:short-incentives} solve private adoption and government ranking. Section~\ref{sec:short-stability} examines institutional stability and robustness. Section~\ref{sec:short-discussion} discusses interpretation and limits.

\section{Model and product-market payoffs}\label{sec:short-model}
Three countries $N=\{1,2,3\}$ each have one firm selling in all three segmented, unit-size national markets; quantities in destination $k$ are $q_i^k$. Governments choose a formal standards partition $\rho$: separate standards $SW=\{\{1\},\{2\},\{3\}\}$, one of three two-country regional unions $SU_{ij}=\{\{i,j\},\{o\}\}$, or international standardization $IS=\{N\}$. A firm natively supports its own coalition's standard. It can pay real fixed cost $F>0$ \emph{per additional standard}, with that adoption covering every country using the additional standard. In each destination the firm uses the destination's standard if it supports it, and its native standard otherwise. The legal/formal partition is unchanged by private adoption.

In each market consumers have heterogeneous valuations $u\sim U[0,1]$, buy at most one unit, and obtain $u+g_i^k-p_i^k$ from firm $i$. For a compatibility group $G_i^k$ of firms using the same standard, let $g_i^k=v\sum_{j\in G_i^k}q_j^k$ if $|G_i^k|\ge2$, and $g_i^k=0$ for a singleton. With total quantity $Q_k=\sum_jq_j^k$, the interior inverse demand is
\begin{equation}\label{eq:short-demand}
p_i^k=1-Q_k+g_i^k,\qquad 0<v<\tfrac14.
\end{equation}
The singleton rule treats network benefits as gains from interoperability across producers: a proprietary product alone generates none. It does not model a proprietary installed base. This assumption materially affects government rankings; allowing singleton benefits eliminates the initial regional advantage on the same domain (Online Supplement E).

The intrinsic production cost is zero. Selling without supporting the destination's standard entails marginal cost $c>0$; native support is costless. A firm that newly adopts a standard obtains full compatibility but may retain residual marginal cost $d=\lambda c$, $0\le\lambda\le1$. Complete relief is $d=0$, whereas $\lambda=1$ still changes compatibility even without marginal-cost relief. We impose
\begin{equation}\label{eq:short-domain}
0<c<\frac{1-3v}{3(1-v)},\quad F>0,
\end{equation}
which ensures interior equilibria of the canonical Cournot configurations. Firms choose additional standards simultaneously after $\rho$ is specified, then compete in quantities; all product-market and adoption stages are re-solved for each candidate formal partition.

For market mass $m_k$ (unity in the symmetric benchmark), the worldwide profit of home firm $i$ is $\Pi_i=\sum_k m_k(p_i^k-r_i^k)q_i^k-Fn_i$, with $r_i^k\in\{0,c,d\}$ its actual destination cost and $n_i$ the number of additional standards. National welfare is $W_i=CS_i+\Pi_i$; in this linear unit-demand environment $CS_k=Q_k^2/2$ (times $m_k$ when markets differ in size). Foreign producer profit is excluded from domestic national welfare. No transfers are allowed.

Let $P$ and $K_I$ be the per-firm operating profit and member consumer surplus under $IS$. Under a symmetric $SU$, before adoption, $A$ and $B$ are respectively a member and the outsider's profits in a member market; $C$ and $D$ are respectively a member's and the native outsider's profits in the outsider market. Write $K_M$ for a member market's consumer surplus and $K_O$ for outsider-market surplus. Direct Cournot first-order conditions give
\begin{align}\label{eq:short-blocks}
P&=\frac1{16(1-v)},& K_I&=\frac9{32(1-v)^2},\nonumber\\
A&=\frac{(1-v)(1+c)^2}{4(2-3v)^2},&B&=\frac{[1-3c-3(1-c)v]^2}{4(2-3v)^2},\nonumber\\
C&=\frac{(1-v)(1-2c)^2}{4(2-3v)^2},&D&=\frac{(1+2c-3v)^2}{4(2-3v)^2},\nonumber\\
K_M&=\frac{(3-c-3v+3cv)^2}{8(2-3v)^2},&K_O&=\frac{(3-2c-3v)^2}{8(2-3v)^2}.
\end{align}
Under $SW$, the market-level blocks are
\begin{equation}\label{eq:short-SW-blocks}
H=\frac{(1+2c)^2}{16},\qquad
S=\frac{(1-2c)^2}{16},\qquad
K_W=\frac{(3-2c)^2}{32}.
\end{equation}
For example, a member-market Cournot equilibrium has two identical member quantities $x=(1+c)/[2(2-3v)]$ and outsider quantity $y=[1-3c-3(1-c)v]/[2(2-3v)]$. Their first-order conditions are $1-3(1-v)x-y=0$ and $1-2x-2y-c=0$, giving $A=(1-v)x^2$, $B=y^2$ and $K_M=(2x+y)^2/2$. This supplies a direct derivation bridge to the government welfare comparisons. Here $H$, $S$, and $K_W$ are the domestic profit, foreign profit, and consumer surplus per market under separate standards $SW$. Thus $W^{SW}=K_W+H+2S$ and $W^{IS}=K_I+3P$. All remaining Cournot configurations and interiority proofs appear in Online Supplement A--B.

\section{Endogenous private adoption}\label{sec:short-adoption}
The strategic adoption margin must be checked before evaluating government payoffs. Under complete marginal-cost relief ($d=0$), let
\begin{equation}\label{eq:short-thresholds}
T_A=P-B,\quad T_U=A-C,\quad T_W=A-S,\quad S=\frac{(1-2c)^2}{16}.
\end{equation}
The outsider gains $2T_A-F$ by supporting a regional union's standard because one investment covers two markets; a member's reverse adoption affects only the outsider's market. Under $SW$, foreign firms face either $T_W-F$ or $T_A-F$. Complete best-response calculations in Online Supplement C imply the following sufficient selection-free condition:
\begin{equation}\label{eq:short-F-region}
F_L=\max\{T_W,T_A\}<F<2T_A=F^*.
\end{equation}
In this interval the outsider uniquely adopts the bloc standard in every $SU$, members uniquely refrain from reciprocal adoption, and $SW$ generates no private adoption. For $F>F^*$ there is no adoption under either $SU$ or $SW$. These are sufficient statements on the canonical domain, not a priority claim about asymmetric adoption.

For incomplete cost relief, define the outsider's updated per-member-market gain $T_O(d)=B_R(d)-B$, where $B_R(d)$ is solved below. Since reverse adoptions must also be reoptimized, let $T_W(d)$ and $T_A^R(d)$ denote the exact largest unilateral $SW$ adoption gain and the reverse-member gain conditional on another member's adoption, respectively; their analytical expressions and all other deviation cases appear in Online Supplement D. With
\begin{equation}\label{eq:short-F-residual}
F_L(d)=\max\{0,T_W(d),T_A^R(d)\},\quad F^*(d)=2T_O(d),
\end{equation}
the strict interval $F_L(d)<F<F^*(d)$ likewise selects outsider-only $SU$ adoption and no $SW$ adoption. Its validity relies on the finite full adoption-game deviation audit, not on imposing an arbitrary equilibrium selection. A nonempty interval is exhibited in Proposition~\ref{prop:incomplete-adaptation-open-set}.

\section{Private Adaptation and Government Incentives}
\label{sec:short-incentives}

We now turn from the firm's adoption decision to the government's comparison of formal standards regimes. This is the paper's central mechanism. The key question is not simply whether the outsider adopts the bloc standard, but how much of the member-market rent survives after adoption.

Consider a symmetric regional standards union with two members and one outsider. Let $W_M^{SU,N}$ denote a member government's continuation payoff before private bypass and $W^{IS}$ its payoff under international standardization.

\subsection{Pre-adoption decomposition}
\label{subsec:selective-decomposition}

Using the Cournot blocks from Section~\ref{sec:short-model},
\begin{equation}
    W_M^{SU,N}=K_M+2A+C,
    \qquad
    W^{IS}=K_I+3P.
    \label{eq:canonical-welfare-pairs}
\end{equation}
Define the member-market component
\begin{equation}
    \mathscr E
    \equiv
    (K_M-K_I)+2(A-P)
    \label{eq:exclusion-surplus}
\end{equation}
and the reciprocal outsider-market disadvantage
\begin{equation}
    \mathscr D
    \equiv
    P-C.
    \label{eq:reciprocal-disadvantage}
\end{equation}
Then
\begin{equation}
    \boxed{
    W_M^{SU,N}-W^{IS}
    =\mathscr E-\mathscr D.
    }
    \label{eq:selective-before}
\end{equation}
The labels are descriptive rather than sign assumptions. $\mathscr E$ is a net national-welfare effect combining member-market consumer surplus and the domestic firm's profits in both member markets. $\mathscr D$ captures the domestic firm's profit disadvantage in the outsider market relative to IS.

For the political reversal studied below, the relevant initial configuration satisfies
\begin{equation}
    \mathscr E>\mathscr D>0.
    \label{eq:initial-preference-condition}
\end{equation}
The first inequality means the member initially prefers the regional union; the second ensures a genuine reciprocal disadvantage remains outside the bloc.

\subsection{Incomplete private adaptation}
\label{subsec:partial-adaptation}

The full-bypass baseline assumes that adopting the bloc standard both creates full compatibility and eliminates the outsider's marginal adaptation cost in member markets. To separate those two effects, let adoption create full compatibility but leave residual marginal cost
\begin{equation}
    d=\lambda c,
    \qquad
    0\leq\lambda\leq1.
    \label{eq:residual-cost-definition}
\end{equation}
Thus $\lambda=0$ reproduces complete cost relief, $0<\lambda<1$ represents incomplete relief, and $\lambda=1$ changes compatibility without reducing the original marginal adaptation burden. The endpoint $\lambda=1$ is therefore not the no-adoption state.

After outsider adoption, all three firms use the bloc standard in a member market. The two member firms have zero marginal cost and the outsider has cost $d$. Re-solving the Cournot subgame gives
\begin{equation}
    q_R^M(d)=\frac{1+d}{4(1-v)},
    \qquad
    q_R^O(d)=\frac{1-3d}{4(1-v)}.
    \label{eq:partial-erosion-quantities}
\end{equation}
The corresponding member-firm profit, outsider profit, and consumer surplus are
\begin{align}
A_R(d)&=\frac{(1+d)^2}{16(1-v)},&
B_R(d)&=\frac{(1-3d)^2}{16(1-v)},\nonumber\\
    K_R(d)&=\frac{(3-d)^2}{32(1-v)^2}.
    \label{eq:partial-erosion-blocks}
\end{align}
Because $0\leq d\leq c<1/3$, the outsider remains active throughout the certified canonical domain.

A regional member has not adopted the outsider's standard, so its outsider-market profit remains $C$. Hence
\begin{equation}
    W_M^{SU,O}(d)
    =
    K_R(d)+2A_R(d)+C.
    \label{eq:partial-erosion-welfare}
\end{equation}
Subtracting $W^{IS}$ produces the exact identity
\begin{equation}
    \boxed{
    W_M^{SU,O}(d)-W^{IS}
    =
    -\mathscr D+R(d,v),
    }
    \label{eq:partial-erosion-identity}
\end{equation}
where
\begin{equation}
    \boxed{
    R(d,v)
    =
    \frac{d\bigl[(5-4v)d+2(1-4v)\bigr]}
         {32(1-v)^2}.
    }
    \label{eq:residual-rent}
\end{equation}
The residual rent combines lower home consumer surplus with higher domestic operating profit relative to IS:
\[
R(d,v)=\underbrace{-\frac{d(6-d)}{32(1-v)^2}}_{K_R-K_I}
       +\underbrace{\frac{d(2+d)}{8(1-v)}}_{2(A_R-P)}.
\]
The outsider's remaining cost protects member firms in both member markets. Its removal erodes that protection while leaving their outsider-market disadvantage $\mathscr D$ unchanged. Thus $R$ is an equilibrium welfare component, not an imposed erosion parameter.

On $0\leq v<1/4$ and $d\geq0$,
\begin{equation}
    R(0,v)=0,
    \qquad
    \frac{\partial R}{\partial d}
    =
    \frac{2(5-4v)d+2(1-4v)}
         {32(1-v)^2}>0.
    \label{eq:residual-rent-monotonicity}
\end{equation}
The amount of residual rent therefore rises monotonically with the remaining adaptation cost.

\subsection{Success, intermediate, and failure regions}
\label{subsec:residual-rent-regions}

The government's response follows directly from equations~\eqref{eq:selective-before} and \eqref{eq:partial-erosion-identity}.

\begin{proposition}[Residual-rent government effect]
\label{thm:selective-erosion}
Suppose $\mathscr E>\mathscr D>0$. Then private outsider adaptation with residual cost $d$ has three distinct effects:
\begin{enumerate}
    \item if $R(d,v)<\mathscr D$, the member's ranking reverses from SU before adoption to IS after adoption;
    \item if $\mathscr D<R(d,v)<\mathscr E$, the relative incentive for IS increases, but SU remains preferred;
    \item if $R(d,v)>\mathscr E$, private adaptation weakens the relative incentive for IS.
\end{enumerate}
At $R(d,v)=\mathscr D$, the member is indifferent between SU and IS after adoption. At $R(d,v)=\mathscr E$, adoption leaves its relative IS incentive unchanged.
\end{proposition}

\begin{proof}
Before adoption the SU-minus-IS gap is $\mathscr E-\mathscr D>0$. After adoption it is $-\mathscr D+R(d,v)$, so the post-adoption ranking favors IS exactly when $R(d,v)<\mathscr D$. The change in the relative incentive for IS is
\[
\bigl(W^{IS}-W_M^{SU,O}(d)\bigr)
-
\bigl(W^{IS}-W_M^{SU,N}\bigr)
=
\mathscr E-R(d,v),
\]
which is positive exactly when $R(d,v)<\mathscr E$. Since $\mathscr E>\mathscr D$, the three regions follow.
\end{proof}

For any target $Z>0$, strict monotonicity implies a unique positive solution to $R(d,v)=Z$:
\begin{equation}
    d_Z
    =
    \frac{
      -(1-4v)+
      \sqrt{(1-4v)^2+32(5-4v)(1-v)^2Z}
    }{5-4v}.
    \label{eq:residual-cost-threshold}
\end{equation}
Let $d_D=d_{\mathscr D}$ and $d_E=d_{\mathscr E}$. Under \eqref{eq:initial-preference-condition},
\[
0<d_D<d_E.
\]
Thus the success, intermediate, and failure regions can equivalently be read from the position of $d$ relative to $d_D$ and $d_E$, truncated to the economically admissible interval $0\leq d\leq c$.

Table~\ref{tab:residual-rent-regions} summarizes the regions and their boundaries.

Figure~\ref{fig:selective-erosion} illustrates this logic. It is a numerical visualization of the analytical thresholds, not a proof.
\begin{figure}[H]
    \centering
    \includegraphics[width=\linewidth]{Fig1.pdf}
    \caption{Residual rent and government incentives. The solid curve $R(d,v)$ is compared with the dashed threshold $\mathscr D$ and dotted threshold $\mathscr E$; vertical lines mark $d_{\mathscr D}$ and $d_{\mathscr E}$. At $(c,v)=(0.073,0.225)$, the three regions show ranking reversal, incentive strengthening without reversal, and incentive weakening. This is a numerical illustration of Proposition~\ref{thm:selective-erosion}, not an adoption-equilibrium certificate}
    \label{fig:selective-erosion}
\end{figure}

\subsection{The firm continuation must hold at the same primitives}
\label{subsec:joint-adoption-government}

A government-ranking condition alone is not enough. The outsider-only adoption state must itself be selected by firms after residual costs are introduced. Let
\begin{equation}
    T_O(d)=B_R(d)-B
    \label{eq:partial-outsider-threshold}
\end{equation}
denote the outsider's per-member-market gain from adopting the bloc standard. Re-solving reverse adoption under residual cost gives thresholds $T_U(d)$, $T_W(d)$, and $T_A^R(d)$ for unilateral SU-member adoption, unilateral SW adoption, and adoption after the other member has already adopted. Their closed forms are reported in Online Supplement, Appendix D.

Define
\begin{equation}
    F_L(d)
    =
    \max\{0,T_W(d),T_A^R(d)\},
    \qquad
    F^*(d)=2T_O(d).
    \label{eq:partial-fixed-cost-region}
\end{equation}
If
\begin{equation}
    F_L(d)<F<F^*(d),
    \label{eq:partial-selection-free-region}
\end{equation}
the outsider strictly adopts the bloc standard, SU members strictly do not adopt the outsider standard against either member action, and SW has no private adoption. The continuation is therefore selection free.

Combining \eqref{eq:partial-selection-free-region} with Proposition~\ref{thm:selective-erosion} gives the economically relevant joint region.

\begin{proposition}[Incomplete adaptation with actual preference reversal]
\label{prop:incomplete-adaptation-open-set}
There exists a nonempty open set of primitives with $0<\lambda<1$ such that: (i) a regional member initially prefers SU to IS; (ii) the outsider-only SU adoption profile is selection free; (iii) SW induces no private adoption; and (iv) after outsider adoption the regional member strictly prefers IS.
\end{proposition}

\begin{proof}
At
\[
(c,v,\lambda,F)
=
\left(
\frac1{10},
\frac6{25},
\frac12,
\frac9{100}
\right),
\qquad
d=\frac1{20},
\]
the pre-adoption gap is
\[
W_M^{SU,N}-W^{IS}
=
\frac{2408509}{295731200}>0,
\]
while the post-adoption gap is
\[
W_M^{SU,O}(d)-W^{IS}
=
-\frac{1341}{184832}<0.
\]
The re-solved fixed-cost thresholds satisfy
\[
F_L(d)
=
\frac{2122041}{31129600}
<
\frac9{100}
<
\frac{459189}{3891200}
=
F^*(d).
\]
All inequalities are strict and all relevant expressions are continuous on the interior domain. A nonempty open neighborhood therefore preserves the same adoption pattern and ranking reversal.
\end{proof}

The proposition rules out the interpretation that the government effect is an artifact of exact cost elimination. It does not establish portability across demand systems or competition modes. Section~\ref{sec:short-stability} reports the tests that limit that broader interpretation.


\subsection{Conduct-neutral mechanism: market access and national incidence}
\label{subsec:general-mechanism}
The preceding propositions concern the certified Cournot environment and
remain unchanged.  To identify which economics travels beyond that
environment, consider the same two-member/one-outsider standards game but
leave the product-market conduct and demand functions unspecified.  In a
representative member destination, let \(q=(q_1,q_2,q_O)\),
\(p(q,a)\) be differentiable inverse prices, and \(S(q,a)\) be its
welfare-consistent consumer-surplus function.  We do \emph{not} require
\(p=\nabla U\) or impose a particular market conduct.  Demand,
prices, and surplus must nevertheless arise from the same consumer
optimization problem with fulfilled network expectations, as specified
below; the canonical convention in Section~\ref{sec:short-model}
is an exact member of this economically admissible class.
The continuously varying \(a\in[0,1]\) is a counterfactual
interpolation from no private adoption to outsider-only adoption, not
a third discrete strategy available to firms.  Assume that along this
path the product-market equilibrium \(q(a)\) is regular and
differentiable, home production costs \(C_h(q_h)\) are differentiable,
the two member destinations are symmetric, and the home firm's
operating profit in the outsider destination does not change.



\paragraph{Consumer microfoundation and fulfilled network expectations.}
In each destination a continuum of atomistic consumers has unit
demands and quasi-linear numeraire utility.  A buyer with taste vector
\(t\) obtains \(u_i(t,a)=\theta(t)+\varepsilon_i(t)+g_i^e-p_i\)
from standard-compatible product \(i\), or zero from the outside
option.  A buyer takes expected network benefits \(g_i^e\) as given;
fulfilled expectations require \(g_i^e=G_i(q,a)\) at the market
equilibrium.  The firm, by contrast, may internalize
\(G_i(q,a)\) in its chosen quantity or price strategy.
The resulting national welfare is defined from \emph{realized
buyer utility net of prices}, not from an arbitrary pair of
inverse demands and postulated consumer-surplus functions.

The original specification is nested \emph{exactly}, not as a
limit of an unrelated differentiated-demand system.  Set
\(\varepsilon_i(t)=0\), let \(\theta\) be uniformly distributed
on \([0,1]\), and let active brands have common
network-adjusted price \(z=p_i-g_i^e\).  Buyers with
\(\theta\ge z\) purchase, with the quantities supplied by
the Cournot firms clearing the allocation among indifferent
buyers.  Then
\[
 Q=1-z,\qquad p_i=1-Q+G_i(q,a),\qquad
 CS=\int_{1-Q}^{1}\bigl(\theta-(1-Q)\bigr)d\theta
 =\frac{Q^2}{2}.
\]
For any continuously differentiable, strictly decreasing inverse
quantile \(D(Q)\) of a general willingness-to-pay distribution,
the same consumer optimization instead gives
\begin{equation}
 p_i=D(Q)+G_i(q,a),\qquad
 CS=\int_0^Q D(z)\,dz-QD(Q).
 \label{eq:general-consumer-foundation}
\end{equation}
The original uniform case is \(D(Q)=1-Q\); its
group-specific, singleton-zero \(G_i\) and Cournot quantities
are \emph{unchanged}.  With nondegenerate brand tastes
\(\varepsilon_i\), consumers need not be indifferent,
and differentiable product-specific demands can support
regular price-setting equilibria, but their existence and
the sign of government preferences remain separate questions.

For a no-tie regular buyer-allocation branch, denote by
\[
 V(p,g)=\int\max\{0,\max_i[
   \theta(t)+\varepsilon_i(t)+g_i-p_i]\}\,d\mu(t)
\]
the quasilinear consumer surplus.  Buyer optimization gives
\(V_{p_i}=-q_i\) and \(V_{g_i}=q_i\).
Along a differentiable, fulfilled-expectations market path,
\(S(q(a),a)=V(p(a),G(q(a),a))\), so the total differential is
\begin{equation}
 \frac{dS}{da}
 =\sum_{i=1,2,O}q_i(a)
   \left[\frac{dG_i}{da}-\frac{dp_i}{da}\right].
 \label{eq:microfounded-consumer-envelope}
\end{equation}
For the tied homogeneous-buyers Cournot case the identical
pathwise formula follows directly from
equation~\eqref{eq:general-consumer-foundation} and
\(p_i=D(Q)+G_i\), without applying a differentiable
product-choice envelope at a point of mass indifference.

Combining this \emph{microfounded} surplus differential with
national producer-profit accounting gives the more revealing
specialization of equation~\eqref{eq:general-incidence}:
\begin{equation}
 \boxed{
 \frac{dW_M^{SU}}{da}
 =\underbrace{\sum_i q_i\frac{dG_i}{da}}_{
   \text{buyer network benefits}}
 -\underbrace{q_O\frac{dp_O}{da}}_{
   \text{foreign-price incidence}}
 +\underbrace{\sum_{h=1}^{2}(p_h-C'_h)
                         \frac{dq_h}{da}}_{
   \text{home markup displacement}}.}
 \label{eq:microfounded-national-incidence}
\end{equation}
All derivatives include the total equilibrium response to
the hypothetical adoption path.  In particular, the
buyer network term is \emph{not} dropped merely because
the original model happens to have \(CS=Q^2/2\):
the benefit is capitalized into fulfilled-expectations
prices.  This identity does not use the Cournot or Bertrand
first-order conditions and is consistent with the
model-independent private-investment and integrated-welfare
sufficiency tests below.  Its individual incidence components
are familiar welfare accounting, not independently new claims.

National accounting then gives
\[
 W_M^{SU}(a)=S(q(a),a)+
 \sum_{h=1}^{2}\{p_h(q(a),a)q_h(a)-C_h(q_h(a))\}
 +\pi_M^{O}.
\]
The sum of the \emph{two} member firms' profits in one member
destination equals one home firm's total profits across the two
symmetric member destinations; it does not attribute foreign-owned
profits to a home government.  Put
\begin{align*}
 H_j&=S_{q_j}+\sum_{h=1}^{2}q_h p_{h,q_j}
       +\mathbf1_{\{j=1,2\}}(p_j-C'_j),\\
 A&=S_a+\sum_{h=1}^{2}q_h p_{h,a}.
\end{align*}
These marginal national-incidence coefficients are calculated
from demand, consumer surplus, production costs and market responses,
not from the sign of the \emph{ex post} government ranking.  By the
chain rule, for either quantity- or price-competition equilibria,
\begin{equation}
 \frac{dW_M^{SU}}{da}=A+\sum_{j=1,2,O}H_jq'_j(a).
 \label{eq:general-incidence}
\end{equation}
Where outsider sales increase, \(q'_O>0\), define the domestic-sales
diversion ratios \(\delta_h=-q'_h/q'_O\).  Hence the
\emph{exact, conduct-neutral} marginal national-welfare loss per unit
of expanded outsider sales is
\begin{equation}
 L(a)=\sum_{h=1}^{2}H_h(a)\delta_h(a)-H_O(a)
       -\frac{A(a)}{q'_O(a)},\qquad
 \frac{dW_M^{SU}}{da}=-q'_O(a)L(a).
 \label{eq:general-displacement}
\end{equation}
\(H_h\delta_h\) is domestic producer displacement, including induced
prices and consumer-surplus effects; \(H_O\) captures the outsider
quantity incidence; \(A\) records direct compatibility and home-price
effects.  Their signs are \emph{not} assumed universal.

\begin{proposition}[Conduct-neutral sufficient conditions for reversal]
\label{prop:general-conduct-neutral}
Suppose the same fixed-cost adoption game is solved under \(SU\) and
\(SW\).  Let \(g_O^0\) be the outsider's gross profit gain from adding
the member standard when both member firms do not adopt;
\(\bar g_H\) the maximum gross gain of any member from reverse
adoption over \emph{all} other \(SU\) adoption profiles; and
\(\bar g_{SW}\) the maximum incremental gross gain from adding one
foreign standard over \emph{all} \(SW\) adoption profiles.
If
\[
 \max\{0,\bar g_H,\bar g_{SW}\}<F<g_O^0,
 \tag{G1}
\]
the unique firm-stage equilibrium has outsider-only adoption under
\(SU\) and no adoption under \(SW\).  In addition, suppose that the
regular path above satisfies \(q'_O>0,\ \delta_h\ge\underline\delta_h
\ge0,\ H_h\ge\underline H_h\ge0,\ H_O\le\overline H_O,\)
and \(A/q'_O\le\overline A\) everywhere.  Write
\[
 \kappa=\sum_{h=1}^{2}\underline H_h\underline\delta_h
       -\overline H_O-\overline A,\quad
 \Delta q_O=q_O(1)-q_O(0)>0.
\]
If the initial counterfactual member advantage satisfies
\begin{equation}
 0<G_0\equiv W_M^{SU}(0)-W^{IS}
   <\kappa\,\Delta q_O,
 \tag{G2}
\end{equation}
then private outsider adoption reverses the member government's
ranking from \(SU\) to \(IS\), without restrictions to linear demand
or Cournot conduct.  Strict adoption and ranking margins persist
under sufficiently small regular payoff perturbations.
\end{proposition}
\begin{proof}
The first inequality in (G1) strictly dominates each member's
reverse adoption under every rival profile, and each incremental
\(SW\) foreign-standard addition; after eliminating these actions,
the outsider strictly gains from adoption by the second inequality.
Thus the stated firm continuation is unique.  Under the response
bounds, equation~\eqref{eq:general-displacement} yields
\(L(a)\ge\kappa\).  Integration gives
\(W_M^{SU}(1)-W_M^{SU}(0)\le-\kappa\Delta q_O<-G_0\).
The post-adoption member therefore prefers \(IS\).  Continuity
preserves the finite strict inequalities wherever equilibria remain
well-defined and regular.
\end{proof}

The proposition supplies \emph{checkable sufficient inequalities},
not an assumption-free prediction that private interoperability
always favors international harmonization.  The original model's
less restrictive adoption interval and explicit \(R(d,v)\) conditions
remain the controlling sharper results for its own primitives.
Also, the pointwise lower bound is sufficient, not necessary:
national incidence can change sign along the interpolation while
its integrated effect still reverses the ranking.  In particular,
the differentiated-demand failure at \(\gamma=1/2\) in Online
Supplement E remains a genuine failure, not a contradiction.


\paragraph{When local welfare effects change sign.}
The preceding pointwise condition is deliberately sufficient but can be
too demanding even in the original certified Cournot model.  Write
\(I(a)=A(a)+\sum_jH_j(a)q'_j(a)\), using the
conduct-neutral incidence coefficients of
equation~\eqref{eq:general-incidence}.  The following alternative admits
a sign change in the marginal national-welfare effect.

\begin{corollary}[Integrated incidence with non-monotone welfare]
\label{cor:integrated-incidence}
If \(I(a)\) is nondecreasing on \([0,1]\) and, for some integer
\(n\ge1\),
\begin{equation}
  \frac1n\sum_{\ell=1}^{n}I(\ell/n)<-G_0,
  \label{eq:general-integrated-certificate}
\end{equation}
then outsider-only private adaptation reverses the member's
initial \(SU\) preference whenever the adoption continuation is
selected as specified above.
\end{corollary}
\begin{proof}
On each cell \([(\ell-1)/n,\ell/n]\), monotonicity gives
\(I(a)\le I(\ell/n)\).  Consequently
\(W_M^{SU}(1)-W_M^{SU}(0)=\int_0^1 I(a)\,da\le
n^{-1}\sum_{\ell=1}^n I(\ell/n)<-G_0\).
\end{proof}
For the original Cournot witness
\((c,v,d)=(1/10,6/25,1/20)\), a symmetric partial-compatibility
interpolation with proportional residual-cost relief connects
\emph{exactly} the original pre- and post-adoption market equilibria.
Its \(I(a)\) is strictly increasing by an exact polynomial positivity
certificate, yet \(I(0)<0<I(1)\): the local welfare effect reverses
sign during adaptation.  Nevertheless, the \(n=16\) exact upper sum is
approximately \(-0.0096983\), below
\(-G_0\approx-0.0081443\), so the general sufficient test certifies
the original reversal without requiring welfare to decrease at
every intermediate step.  The interpolation is a proof device, not
a change in the discrete firm-adoption game.  The exact expressions
and sign certificates are reported in the reproducibility audit.
This finite test establishes \emph{applicability}, not independent
novelty of integration or pass-through accounting.

\section{Coalition stability and scope}\label{sec:short-stability}
We consider the static five-partition game. A deviation by coalition $D\subseteq N$ strictly blocks partition $\rho$ if every $i\in D$ obtains higher continuation welfare in the induced $\rho'$, after firms reoptimize standard adoption. Let $\mathcal S_S(F)$ be the unblocked partitions. The full-bypass institutional application below sets $d=0$, and is distinct from the positive-$d$ result in Proposition~\ref{prop:incomplete-adaptation-open-set}. The admissible deviating set and exclusivity rule are the same as in the full canonical partition game; alternative coalitions cannot impose transfers or ask a nonmember to sacrifice national welfare.

On the nonempty canonical region $\Omega_0$ where $\mathscr E>\mathscr D>0$, $P>C$, and $F_L<F^*=2T_A$, full outsider bypass yields a member's welfare $W_M^{SU,O}=K_I+2P+C=W^{IS}-(P-C)$. For the outsider, $W_O^{SU,O}=K_O+2P+D-F$, so $W^{IS}-W_O^{SU,O}=J+F$, where $J=K_I+P-K_O-D>0$ on the domain (Online Supplement B). Thus all three countries strictly prefer $IS$ to an outsider-only $SU$ continuation.

\begin{theorem}[Full-bypass stability application]\label{thm:short-stability}
For $(c,v)\in\Omega_0$, under strict blocking: (i) for $F>2T_A$, $\mathcal S_S(F)=\{SU_{12},SU_{13},SU_{23}\}$; and (ii) for $F_L<F<2T_A$, $\mathcal S_S(F)=\{IS\}$. The high-cost symmetric stable set depends on strict rather than weak/Pareto blocking. The intermediate unique-$IS$ result survives weak/Pareto blocking and a sufficiently small market-size asymmetry.
\end{theorem}
\noindent\emph{Proof sketch.} In (i), firms do not adopt; a member prefers $SU$ to $IS$ and $SW$, but switching to another $SU$ makes a common member indifferent, excluding a strict deviation. In (ii), every $SU$ selects outsider-only adoption; the member gap $P-C>0$, outsider gap $J+F>0$, and the comparison with $SW$ (Online Supplement E) make $IS$ uniquely unblocked. Weak/Pareto blocking permits a symmetric partner-switching deviation in (i), while strict inequalities sustain (ii). For $m_1=m_2=1$, $m_3=1-\delta$, the high-cost partner payoff difference is $\delta(A-C)>0$; full stability-case and small-asymmetry calculations are in Online Supplement E--F.

\paragraph{Limits of the political result.} First, setting $v=0$ gives $W_M^{SU,N}-W^{IS}=c(13c-6)/32<0$ for $0<c<1/3$: a regionally favorable initial ranking is absent, although private one-way adoption can still occur. Second, allowing singleton compatibility groups to enjoy their own network effects also eliminates the initial regional advantage throughout the tested canonical parameter domain; this is a particular alternative network microfoundation, not a universal absence result. Third, in a pre-specified differentiated-demand setting with substitution $\gamma=1/2$, the initial regional member advantage and member-market exclusion component are negative throughout the audited parameter box under both differentiated Cournot and Bertrand competition, despite feasible outsider-only private adoption. These are substantive portability failures, not merely imperfect proofs. The differentiated-demand and singleton tests do not establish that every other specification fails; full calculations and parameter-domain certificates are provided in Online Supplement E.

\section{Discussion and conclusion}\label{sec:short-discussion}
The firm-stage and political mechanisms should not be conflated. Supporting an additional standard has one-time cost $F$ and may leave residual marginal cost $d$; where outsiders adopt unilaterally, consumers in member markets can benefit while protected domestic producers lose profits. Governments count their own consumers and producers, not the foreign adopter. The residual rent $R(d,v)$ therefore changes incentives for \emph{formal} harmonization without changing the formal partition itself. Conditional on a fixed outsider-only continuation and the canonical market structure, lowering $d$ reduces the member's relative preference for $SU$; it is not an unconditional prediction across technologies or countries.

Our results do not imply that international standardization is the globally first-best policy, nor that the positive-$d$ political ranking determines a stable coalition structure. The full-bypass welfare and blocking comparisons concern only the five formal partitions under specified rules. The distinction between technology and policy requires care. A regulator may observe an increase in products that technically comply with multiple standards and conclude that regional arrangements are becoming redundant; the model shows that the national-welfare incentive to negotiate a broader agreement also depends on the displaced domestic operating profits and on disadvantages that remain in outside destinations. As a result, observing private standard adoption alone cannot identify the sign of the government preference effect. In the canonical ranking-reversal region, this is a strictly national-welfare rather than a global-planner comparison. The model uses segmented linear demand, one firm per country, a singleton-zero network benefit and binary adoption; the identified preference reversal is not portable even to the pre-specified differentiated-demand environment. Empirical use would require separate measures of adoption setup costs, remaining per-unit compliance costs, demand-side network benefits, home producer rents, and whether interoperability is reciprocal. No empirically identified policy intervention is claimed. This conditionality is itself the economic lesson: private adaptation can diminish technological incompatibility without making broader public harmonization either automatically more or automatically less attractive.

\section*{AI-assisted methods and reproducibility}
ChatGPT (OpenAI) assisted with theoretical exploration, algebraic checks,
literature discovery, manuscript editing, and the preparation of verification
code. AI output was not used as a substitute for a mathematical proof or as
an authoritative scholarly reference. The author retains responsibility for
the model assumptions, economic interpretation, theorem statements, and
manuscript. Exact symbolic checks, independent Cournot re-derivations, and
formal verification of selected logical claims are documented in the
anonymous replication package. No empirical data are used.

\section*{Online Supplement (Online Resource 1)}
The anonymous Online Resource 1 provides Cournot derivations, firm-stage deviation checks, coalition-stability comparisons, and exact sign certificates (Appendices A--G). Appendix E identifies the alternative equilibria and their polynomial sign proofs; Appendix G supplies the five-partition deviation ledger. The replication package includes the R2--R9 checks, independent matrix reconstruction, and targeted Lean sources. The short presentation preserves the certified results and the full research manuscript.
\bibliographystyle{elsarticle-harv}
\bibliography{references}
\input{table_residual_rent_regions}
\end{document}
